% Generated draft: requires mathematical and source review.
\documentclass[11pt]{article}
\usepackage[margin=25mm]{geometry}
\usepackage{fontspec}
\setmainfont{Latin Modern Roman}
\usepackage{amsmath,amssymb,mathtools,unicode-math}
\setmathfont{Latin Modern Math}
\usepackage{enumitem,needspace,xurl,hyperref}
\setlength{\parindent}{0pt}
\setlength{\parskip}{5pt}
\setlength{\emergencystretch}{3em}
\newcommand{\N}{\mathbb{N}}
\newcommand{\Z}{\mathbb{Z}}
\newcommand{\Q}{\mathbb{Q}}
\newcommand{\R}{\mathbb{R}}
\newcommand{\C}{\mathbb{C}}
\newcommand{\F}{\mathbb{F}}
\newcommand{\A}{\mathbb{A}}
\newcommand{\PP}{\mathbb{P}}
\newcommand{\E}{\mathbb{E}}
\newcommand{\eps}{\varepsilon}
\newcommand{\dd}{\,\mathrm{d}}
\begin{document}
\section*{UnsolvedMath ID: 10400156}
\textbf{Source problem code:} catalogue entry
\subsection{Definitions and mathematical statement}
Let \((V,Z)\) be the given \((2+1)\)-dimensional TQFT in the source's conventions, and let \(\Sigma_g\) be a closed oriented surface of genus \(g\). Its mapping class group is \(\mathcal M_g=\pi_0(\operatorname{Diff}^+(\Sigma_g))\). The source writes \(\widetilde{\mathcal M}_g\), presumably an extension or enhanced version needed to define the TQFT action; that group is not defined in the extract. Denote the induced action by
\[
 \rho_g:\widetilde{\mathcal M}_g
       \longrightarrow\operatorname{End}(V(\Sigma_g)).
\]
Maps arising from invertible mapping cylinders are invertible, so the image lies in the automorphism group when the usual functorial action applies.

Problem 8.10 asks, for the given theory, to determine whether
\[
 \left|\rho_g(\widetilde{\mathcal M}_g)\right|<\infty.
\]
It is a request to decide or characterize finiteness of the image, not a universal assertion that all such images are finite. The intended genus range, the TQFT parameters and coefficients, and the precise extension must be specified. Finiteness of this actual image is not interchangeable with finiteness after projectivizing the representation: central scalar images can affect the answer. The record gives no finite-image criterion, algorithmic input format, or specified family of TQFTs.
\subsection{Short English statement}
For a specified TQFT and genus, determine whether the induced action of the source's extended mapping class group has finite image. The exact extension and theory parameters are missing, and the question concerns the stated linear image rather than an automatically substituted projective image.
\subsection{Sources}
\begin{itemize}
\item Local UnsolvedMath registry, permanent ID \(10400156\), exact source code \texttt{catalogue entry}, Problem 8.10. No individual proposer or proposed year is supplied.
\item Ohtsuki, \emph{Problems on invariants of knots and 3-manifolds} (2002), Problem 8.10, PDF page 125: \url{https://msp.org/gtm/2002/04/gtm-2002-04-024s.pdf}. Bold printed headings distinguish 231 numbered items from surrounding references.
\item Extraction pdf-font-structure, accessed 2026-07-29. Source evidence NEEDS\_REVIEW presents an open problem and calls for a current primary-literature audit. Supplied research classification OPEN-TRIAGE identifies mapping-class-image finiteness without an answer or further primary references.
\item Dataset attribution: Ulam AI / UnsolvedMath Contributors, CC BY 4.0. Source rights NEEDS\_REVIEW: publicly accessible author/publisher copy, redistribution terms requiring review.
\end{itemize}
\end{document}
